\chapter{Related Work}\label{ch:related-work}

% =====================================================================
% SKELETON ONLY -- structure agreed 2026-09-23 (manifest §6). No prose yet;
% each section is elaborated one at a time.
%
% BUDGET: 4-5 pp ~= 2,000-2,250 words, plus a ~100-word chapter opening.
% Five sections. Retrieval fusion has no section of its own: after §3.4 the
% only related-work content left is that ARQMath participants fused runs,
% which is one unelaborated sentence in §2.1.
%
% THE BOUNDARY AGAINST CHAPTER 3
%   Background says what a thing IS. Related Work says what other people
%   BUILT with it, and where this thesis stands.
%   A background sentence is true if this thesis had never been written;
%   a related-work sentence names a system, a result or a position.
%   => This chapter names systems and their published outcomes. It never
%      explains a mechanism. RRF's formula is §3.4; that the best automatic
%      ARQMath-3 Task 1 run was an RRF ensemble is §2.1.
%
% THE ORDERING CONSTRAINT -- this chapter comes BEFORE Chapter 3
%   No term Chapter 3 defines may appear here: not SLT, OPT, path tuple,
%   prime, hole, nDCG', pool, judge. Write coarse instead --
%   "a tree-based formula representation", "systems built after the
%   assessment was finished", "an automatic relevance labeller".
%
% HOUSE STYLE (CLAUDE.md §3): open each section by saying why it exists;
% motivate before defining; carry numbers with their provenance; close by
% naming the problem the next section inherits. Forward references name a
% QUESTION a later section settles, never a location or an outcome.
%
% Per-paper detail, with the exact figures and our position on each, is in
% manifest §16. Twelve of the papers below still need entries in sources.bib.
% =====================================================================

% TODO chapter opening -- PENDING A DECISION (2026-09-23). Section 2.1 now opens the chapter
% by explaining the field, so a separate opening would say the same thing twice. Abram to
% confirm dropping it. Original plan, ~100 w: why a related-work chapter precedes the
% machinery: this thesis assembles known components in an untried
% arrangement and evaluates it on a collection it did not participate in,
% so what is new is the arrangement and the evaluation, not the parts.
% Name the four sections; say the fourth matters most.


\section{Mathematical information retrieval}\label{sec:rw-mir}

% Rewritten 2026-09-24 (Abram) as the FIELD section, replacing the 2026-09-23 version.
% What it covers: what the field is, the three retrieval tasks it defines, the campaigns
% that defined them, the difficulties common to all of them, and the task split that
% Section 2.2 inherits. It does NOT say what the ARQMath collection is -- that is
% Section 3.1, and the earlier draft duplicated it.
%
% THE WORKING TEST used throughout: every sentence carries an agent or a date -- a field,
% a campaign, a group, a year. A sentence that would be true with nobody having done
% anything belongs to Chapter 3 and was cut.
%
% Cut in this rewrite, deliberately:
%  - the three-stage methods history (conventional -> tree-based -> neural). It was a
%    trailer for Section 2.2 and named Tangent-S and Tangent-CFT a page before that
%    section introduces them.
%  - the ARQMath collection description (Math Stack Exchange corpus, the two tasks,
%    topics posed after the snapshot). Section 3.1 carries all of it.
%  - "The task decides which of the two matters more" -- the chapter's one ranking
%    sentence, banned by her instruction of 2026-09-23 (review E4).
%  - "held in Tokyo in 2013"; the NTCIR-12 subtask breakdown (Abram, 2026-09-24).
%  - the structural/semantic split stated in the opening AND again at the close. It now
%    appears once, in the final paragraph, where it carries the argument.
%
% Added in this rewrite:
%  - a paragraph naming the tasks. The old version never enumerated them; formula search
%    and answer retrieval were only ever defined implicitly inside the closing argument.
%  - a challenges paragraph. Did not exist.
%  - the fusion sentence manifest §6 requires and the old draft was missing: participants
%    merged rankings, and the best automatic answer-retrieval run was such an ensemble.
%
% SOURCING (changed 2026-09-24). The NTCIR history and the challenges rested entirely on
% the Dadure survey, a secondary source carrying three campaign generations. They now
% rest on primary sources by the campaigns' own organisers: aizawa2020ntcirmath for the
% NTCIR editions and the open problems, zanibbi2016ntcir12 for the document-retrieval
% task, zanibbi2025mir for the task framework. Dadure is demoted to the single framing
% citation manifest §16 always described it as.
% ⚠ ALL THREE NEW ENTRIES CARRY TODO FIELDS IN sources.bib AND HAVE NO PDF IN
%   thesis/literature/. Their claims here are pitched at what the abstracts support;
%   verify against the PDFs before submission.
%
% Numbers: Tangent-S 0.540 (formula search) / 0.159 (answer retrieval) and 0.272 (text
% baseline) are the organisers' published figures, manifest §16 and R6. The measure is
% never named -- nDCG' is Section 3.5. The sentence after them states that the two tasks
% are scored on different topic sets, so the pair carries a direction and not a size;
% without it the pair reads as a magnitude claim, which §1c-i rules out.
%
% The closing two-explanation sentence is the section's load-bearing move: the second
% explanation is the one RQ1 tests. Do not compress it back into "whether that belongs to
% the representations or to the task".

This thesis puts known components together in a new arrangement, so what it adds can
only be seen against the problems the field already poses. This section describes those
problems, the campaigns that defined them, and the difficulties they share.

Mathematical information retrieval is the problem of searching documents that contain
mathematical notation. It grew out of digital mathematical libraries, where ordinary
text retrieval treated a formula as a string of characters and could not reach the
structure that carries its meaning~\cite{dadure2024mirreview}. A query may be text, a
formula, or both, which makes it a different problem, not a harder one.

Three tasks recur. Formula search takes one formula as the query and returns formulas
like it. Answer retrieval takes a mathematical question, stated in text and formulas,
and returns passages that answer it. Document retrieval, the earliest campaigns' task,
takes keywords and formulas together and returns paragraphs of scientific
articles~\cite{zanibbi2016ntcir12}. Recent benchmarks add premise selection and
statement retrieval over formal mathematics, and find every model weaker where
reasoning is needed~\cite{ju2025mirb}. A recent survey groups them by one distinction:
whether a system finds mathematical content or answers a question with
it~\cite{zanibbi2025mir}.

Shared campaigns defined them. The NTCIR Math tasks came first, from 2013: a pilot over
a hundred thousand arXiv papers, then two further editions that added English
Wikipedia~\cite{aizawa2020ntcirmath}. The ARQMath labs, at CLEF from 2020 to 2022, moved
the corpus to Mathematics Stack Exchange and set answer retrieval beside formula
search~\cite{zanibbi2020arqmath,mansouri2021arqmath2,mansouri2022arqmath3}. Participants
there routinely merged several rankings into one, and the strongest automatic
answer-retrieval run of the final edition was such an
ensemble~\cite{mansouri2022arqmath3}.

Three difficulties run through all of it. Notation is ambiguous: a small set of symbols
stands for an immense number of concepts, so one expression can mean different things in
different documents. The same idea can be written in several ways, so two formulas that
mean the same thing may look nothing alike. And relevance is costly to establish: the
NTCIR organisers drew their assessors from third-year and graduate students of pure
mathematics, and divided documents into paragraphs to hold down the cost of judging
them. Their account of the campaign opens by naming the absence of large-scale
datasets, carefully designed tasks and established evaluation methods as the obstacle
it set out to remove~\cite{aizawa2020ntcirmath}.

Those difficulties fall unevenly across the tasks. A formula query asks for formulas
like the one given, so structural similarity is most of the problem; a question carries
text and several formulas, and a useful answer need not resemble any of them. The
organisers' baselines show it. Tangent-S, run in both tasks, reached 0.540 at formula
search and 0.159 at answer retrieval, against 0.272 for a plain text
baseline~\cite{mansouri2022arqmath3}. The two tasks are scored on different topics, so
the pair shows the same system is much weaker at answer retrieval. The representations
may not carry the similarity answer retrieval rewards; or they may, and no one has
combined them with text retrieval well enough to reach an answer. No team submitted an
answer-retrieval system built on mathematics alone, so the leaderboards cannot tell them
apart. Section~\ref{sec:rw-formula-systems} turns to the representations.


\section{Formula representation and retrieval systems}\label{sec:rw-formula-systems}

% Rewritten 2026-09-24 (Abram), replacing the 2026-09-23 version. Positional only: what
% each system is, what it was measured on, and what it scored. Mechanism belongs to
% Section 3.2.
%
% Cut in this rewrite, deliberately:
%  - MECHANISM, which the 2026-09-23 version explained in paraphrase and Section 3.2
%    defines properly: Tangent-S "reads a formula two ways, once as the arrangement of
%    symbols on the page, once as the operations the expression performs"; Tangent-CFT
%    "flattens each tree into a sequence of fragments, treats those fragments as words,
%    and trains an embedding model over them". Avoiding the banned words (tree, tuple,
%    SLT, OPT) did not make those sentences positional -- they were still the Chapter 3
%    explanation in other clothes. Do not restore them.
%  - the MathBERT bridge into Section 2.3. The old close pivoted on MathBERT's topic
%    classification and headline generation ("no retrieval feeds that generation and no
%    question is being answered"), which hinged the chapter's turn on a paper doing two
%    things unrelated to retrieval.
%  - "A second line learns the representation rather than building it" and the repeated
%    use of *line* (Abram, 2026-09-24). One use survives, in the opening sentence.
%
% Fixed in this rewrite:
%  - SSEMB'S LEADERBOARD CLAIM IS NO LONGER CITED TO SSEMB. The old version attributed
%    "above every automatic run the lab itself recorded" to li2025ssemb; manifest §16
%    warns explicitly against this, because SSEmb's own quoted TangentCFT2 figures
%    disagree with the overview's (0.6868/0.6211 against 0.694/0.641). The leaderboard
%    claim now cites mansouri2022arqmath3. Keep it that way.
%  - the closing sentence manifest §6 asks for and the old version withheld: what this
%    thesis reimplemented and what it did not, stated once. The old file comment sent C2
%    to Section 1.5 and the no-retriever-training concession to Section 7.2 and left this
%    section silent; the manifest wins, and *where this thesis stands* is what Related
%    Work is for. One positional sentence -- not a contribution claim, not a limitation.
%
% ⚠ THE THIRD PARAGRAPH IS THE CHAPTER'S EXPOSED CLAIM. Four ARQMath participant systems
%   combine formula structure with text for ANSWER retrieval -- MathDowsers (tuples from
%   formula trees indexed beside text terms, best Task 1 participant run twice), Approach
%   Zero with Anserini (operator-tree path tokens interpolated with a BM25 run), Approach
%   Zero with ColBERT (structural matching plus a dense retriever, top across all three
%   ARQMath-3 tasks), and Kassaie et al. 2025 (RRF over math and keyword components).
%   None is cited here: Abram dropped them 2026-09-24 to avoid two new bib entries. The
%   paragraph therefore argues from DPRL alone, and claims only what DPRL supports -- that
%   the GROUPING of several formula representations has not been varied where answers are
%   retrieved. It must never be widened back to "formula and text have never been
%   combined for answer retrieval", which those four systems falsify.
%   ⚠ Section 2.5's first gap still says the formula engines have never been entered into
%   answer retrieval "let alone paired with a text retriever" -- reword with §2.5.
%
% Scope verified per paper: Tangent-S and the graph-contrastive work are NTCIR-12 only;
% SSEmb is ARQMath-3 formula search only. Tangent-S's combination figures (bpref 0.524 /
% 0.513 alone, 0.553 combined) stay held out -- a second measure on a second collection,
% and the claim here is ordinal.
% MathAMR is named nowhere: the DPRL sentence was simplified to "raw text with formulas
% left in as written" (Abram, 2026-09-24). Manifest §6 still asks for DPRL/MathAMR.

Section~\ref{sec:rw-mir} left the representations themselves as the open question.
Tangent-S is where the line starts: it indexes two tree representations of a formula
separately and scores a candidate on both, and Davila and Zanibbi report that using the
two together retrieves better than either alone on the NTCIR-12
collection~\cite{davila2017tangents}. That result, on formula search in 2017, is the
published precedent for treating several representations as complementary. Tangent-CFT
replaced the matching with learned vectors and merged its representations by summing
them into one~\cite{mansouri2019tangentcft}. TangentCFT2 reversed that choice, holding
four representations in four indexes and merging their rankings instead, inside formula
search, where the formula system is the whole system~\cite{mansouri2022dprl}.

MathBERT pre-trains a language model over formulas and the text around
them~\cite{peng2021mathbert}; a graph-contrastive method learns formula embeddings with
no relevance labels at all~\cite{wang2024gcl}. Both report gains over Tangent-CFT, and
both were measured on NTCIR-12, a different collection with different queries and
judgments, so neither figure can be set beside the ARQMath ones. SSEmb is the exception,
and the strongest result of the three: it combines a formula's structure with its
surrounding text and reports 0.734 nDCG$'$@10 at ARQMath-3 formula search, the highest
among the automatic systems on its own recomputation of the
leaderboard~\cite{li2025ssemb}. That figure is not on the lab's measure, so it cannot be
set beside the published nDCG$'$ either.

What these representations have not been asked is how they behave as a group.
TangentCFT2 is the one system that holds several formula representations apart, and it
does so only in formula search, where there is no text to combine them with. Its own
team submitted five answer-retrieval runs at the same lab, and none of them used a
formula index at all: they searched the raw text of the posts, with formulas left in it
as written~\cite{mansouri2022dprl}. So the grouping of formula representations, and the
way that group meets text retrieval, has not been varied and measured where answers are
what is retrieved.

This thesis rebuilds the Tangent-CFT-style stack from scratch and keeps its
representations as separate indexes. Its formula embeddings are trained on the corpus's
own formulas, and nothing in it is trained against relevance labels. What it varies is
how those representations are grouped, so the learned systems above are not what it sets
out to beat. Every system in this section is measured on the ranking it produces. None
asks whether a difference visible at the ranking survives into the answer a language
model writes from what was retrieved. Section~\ref{sec:rw-rag} takes that question up.


\section{Retrieval-augmented generation}\label{sec:rw-rag}

% Rewritten 2026-09-24 (Abram), replacing the 2026-09-24 first version. 590 w -> 371 w,
% against a ~400 w budget. This section was the chapter's largest overrun.
%
% TWO CLAIMS WERE FALSE AND ARE NOW GONE. Found 2026-09-24 by search:
%  (1) "Mathematics is barely present in any of this." Retrieval-augmented generation on
%      mathematics is an active area: an Olympiad-level benchmark with a math-aware
%      retrieval task and a retrieval-augmented problem-solving task (MathNet 2026);
%      RAG against GraphRAG on an undergraduate textbook, measuring retrieval accuracy
%      and answer quality separately over 477 question-answer pairs (Chen et al. 2025);
%      a dual-index framework combining vector retrieval with a knowledge graph for
%      school mathematics (Jiang et al. 2026); multi-embedding maths QA (Confident RAG,
%      2025); and RAG over reasoning traces (Arabzadeh et al. 2026). Only MathNet is
%      cited here; the others are in manifest §16's note.
%  (2) "Retrieval quality and answer quality are two measurements, and the first is a
%      poor proxy for the second." Contested, not settled. Salemi and Zamani hold up
%      (229 citations) and Ghosh et al. 2026 corroborate across five benchmarks, but Guo
%      et al. 2025 report a strong positive relationship on one corpus and a weak one on
%      another, and MathNet reports answer quality highly sensitive to retrieval quality.
%      ⚠ The section now says the question is UNSETTLED, which is a better reason for RQ2
%      to exist: if the link were simply weak, measuring it again would prove little.
%      Do not restore the flat claim.
%
% Cut in this rewrite, deliberately:
%  - THE RAG DEFINITION. The old opening defined retrieval-augmented generation -- search
%    step in front of a language model, passages into the prompt, model writes the answer.
%    Section 3.7 italicises and defines that term, with *generator* and *context*. It is a
%    Chapter 3 term and the ordering constraint forbids it here.
%  - "An answer has to be scored for whether it is true, and for whether it follows from
%    what was retrieved" -- Section 3.7's own sentence pair ("whether the answer is
%    grounded in the context it was given, and whether it is correct").
%  - THE RAGAS GLOSS. "which score a generated answer using a language model in place of a
%    human annotator" is near-verbatim §3.7's definition. Manifest §16: RAGAS is NAMED in
%    §2.3 and DEFINED in §3.7. It is named here and nothing more.
%  - Henkel's three prompt conditions (~90 w). They are about prompt guidance, which
%    supports §6.2's single-prompt decision; manifest §16 already assigns Henkel to §6.2.
%    Only the relevance-against-preference finding stays, which is what §2.3 needs.
%  - "It would be natural to assume the two are connected" -- throat-clearing.
%
% Fixed: Henkel's null is no longer stated as "relevance did not predict preference". The
% paper reports no SIGNIFICANT relationship over 51 questions; §1c-i bans that word, so the
% prose gives the sample size and says what they found. Same error still live in §2.5.
%
% Sourcing notes retained from the 2026-09-24 version:
%  - Salemi & Zamani Table 1: Kendall's tau between retrieval scores and downstream RAG
%    performance, five collections, two retrievers. Human relevance annotations top out
%    near 0.2 and are often ~0. Written here as "no higher than about 0.2" -- true of
%    every cell in that block. "Five widely used benchmark collections", not "question
%    answering": two of the five are not QA (Abram, 2026-09-24).
%  - NOT USED HERE, but belongs in §7.2: the same table shows labels assigned by a language
%    model correlate no better with answer quality than human ones. This thesis reads
%    generation quality against judge labels, so §7.2 should concede it.
%  - ARQMath Task 3: Table 5, printed p. 307 (GPT-3 baseline AR 1.346 / P@1 0.500 against
%    approach0 run1 1.282 / 0.436); reusability passage, printed p. 308. The organisers do
%    NOT say the problem is unsolved -- they propose overlap-based measures -- so the prose
%    says they "note a difficulty", not that they could not solve it.
%  - ⚠ alshammari2026mathnet carries TODO fields in sources.bib and has no PDF in
%    thesis/literature/. Its claim here is pitched at what the abstract supports.
%
% NAMING vs DEFINING (Abram, 2026-09-24). The first paragraph names *retrieval-augmented
% generation* as what Lewis et al. called their arrangement. That is attribution, not a
% definition: §3.7 still defines the term, italicised, with *generator* and *context*, and
% Chapter 2 must never italicise it or explain the machinery. The ordering constraint is
% carved out for it in manifest §6 -- the term that names a section may be used inside that
% section, since the heading puts it in front of the reader either way.
%
% ⚠ Section 2.5's second gap leg still says "On mathematical questions almost nothing has
%   been measured at all". That is the claim (1) above, and it must be reworded with §2.5.

A retriever used to be the last stage of a search system. It is now often the first
stage of a system that writes the answer itself, the arrangement Lewis et al.\ set out
and named retrieval-augmented generation~\cite{lewis2020rag}. So what gets judged is the
answer and not the ranking that produced it, and those are two different measurements.

Whether the first carries into the second is unsettled. Salemi and Zamani scored
retrievers with the usual ranking measures against human relevance annotations, compared
those scores against how well the surrounding system answered, and found rank
correlations no higher than about 0.2 across five widely used benchmark collections,
often close to zero~\cite{salemi2024erag}. That result is hard to check against most
published work, because most of it reports only the end of the pipeline: the generated
answer is compared against a known correct answer, and how well the retrieval step did on
its own is never stated, so an improvement cannot be traced to one part or the
other~\cite{salemi2024erag}. Frameworks such as RAGAS are a recent
attempt to make the generation side measurable at all~\cite{es2024ragas}.

On mathematics the picture is different again, and thinner than it looks. Henkel and
colleagues built a tutoring system for middle-school students over a free prealgebra
textbook, annotated how relevant each retrieved section was, and over fifty-one questions
found no relationship between that relevance and which answers readers
preferred~\cite{henkel2024mathrag}. A recent Olympiad-level benchmark reports the
opposite for problem solving, where answer quality is highly sensitive to how good the
retrieval is~\cite{alshammari2026mathnet}. Both search curated corpora of prepared
material by semantic similarity. Neither varies how formulas are represented in the
retriever, so neither says whether a formula index reaches the answer.

ARQMath ran a generation task of its own in 2022, over the same questions as its answer
retrieval task. A baseline that retrieved nothing and answered from the model alone
scored higher than every participating system~\cite{mansouri2022arqmath3}. The organisers
also note a difficulty with reusing what they collected: the judgments made there cannot
be applied to later systems in the way retrieval judgments can, because the measures
depend on matching a new answer against one already judged~\cite{mansouri2022arqmath3}.
That difficulty is not particular to generation, and
Section~\ref{sec:rw-incomplete-judgments} takes it up.


\section{Evaluation with incomplete relevance judgments}\label{sec:rw-incomplete-judgments}

% Rewritten 2026-09-24 (Abram), replacing the 2026-09-24 first version. 707 w -> ~450 w,
% against a ~450 w budget. This was the chapter's last overrun.
%
% Title chosen over "post-hoc evaluation" because this thesis already uses "post-hoc" in
% the other sense -- the §7.2 row disclosing H1.3 as restated after its data was known.
% Two meanings of one word is the defect the review named.
%
% Cut in this rewrite, deliberately:
%  - MECHANISM, twice. The old version explained how bpref works ("looks only at the
%    documents that were judged and asks whether the relevant ones were ranked above the
%    irrelevant ones") and how condensing works ("measures that simply drop the unjudged
%    documents from a ranking and score what is left"). §3.5 defines the prime variants
%    and names bpref. Both are now characterised in a clause, not explained.
%  - THE THESIS ARGUING ON CLARKE & DIETZ'S BEHALF. The old version explained *why*
%    agreement collapses at the top -- "the systems at the top all put plausible documents
%    first; what separates them is whether those documents really answer the question" --
%    which they do not say. ~45 w of unsourced reasoning, now replaced by Fröbe's measured
%    mechanism. Do not restore it, and do not restore "agreement that holds across a whole
%    field but fails at the top is agreement on the comparisons that did not need making":
%    it is rhetoric, not a finding.
%
% Fixed in this rewrite:
%  - A CITATION DEFECT. The old version said Upadhyay et al. "rebuilt the method as an
%    open-source tool, UMBRELA, and then tested it in public: 77 runs from 19 teams",
%    citing upadhyay2025umbrela for both halves. Those are TWO papers: the toolkit
%    (upadhyay2024umbrelatool) and the TREC 2024 RAG Track study (upadhyay2025umbrela,
%    which is the PDF in thesis/literature/). The toolkit claim was unsourced.
%    ⚠ SUPERSEDED 2026-09-29 (citation audit): upadhyay2024umbrelatool was a metadata
%    stub (author "and others", no identifier, arXiv-only) and has been REMOVED. The
%    sentence now cites upadhyay2025umbrela alone, which is where "77 runs from 19 teams"
%    actually appears; the toolkit paper evaluates on TREC DL 2019-23 and never states it.
%  - One finding added from a paper already cited: the TREC study also found that having
%    humans check the model's labels did not improve agreement over full automation. That
%    is independent support for §1c-i decision 2 and for retiring llm_plus_human, beside
%    Abbasiantaeb's. Worth keeping when this section is cut further.
%
% Added 2026-09-24, both earning their place twice:
%  - froebe2025assessors. Eight automatic assessors on TREC DL and the RAG track: they
%    agree with each other more than with humans, and they overrate systems that use a
%    language model to rank by 9-17 positions across 30-50 systems. This is Clarke &
%    Dietz's "biased toward LLM-based rankers" with a number on it. ⚠ ITS SECOND USE IS
%    §4.4: this thesis's retrieval stack contains no language model, so the bias is
%    measured, quantified, and demonstrably does not apply here -- a far stronger reply
%    than the qualitative one §16 currently plans against Soboroff.
%  - szymanski2025judgelimits (key renamed 2026-09-29; arXiv 2024, published IUI 2025).
%    Subject-matter experts agreed with an automatic judge 68%
%    of the time in dietetics and 64% in mental health. ⚠ THIS IS C3'S THIRD LEG,
%    EVIDENCED: the case that automatic labels are unproven where judging needs
%    specialist knowledge no longer rests on assertion. §2.5's third leg survives on it
%    while legs one and two do not.
%
% Constraints honoured: no Chapter 3 vocabulary -- no "pool" (the process is described
% instead), no "prime", no "judge" as a noun for the model, no kappa, no Kendall's tau by
% name. "Rank correlation" is allowed (Abram, 2026-09-24) and carries Clarke & Dietz's
% figures. Thomas et al.'s spread is given qualitatively because those numbers are Cohen's
% kappa, which means nothing without §3.6.
%
% Sourcing notes:
%  - UMBRELA's weaker agreement is RUN LEVEL vs TOPIC LEVEL, not per-document.
%  - Clarke & Dietz: 0.89 over the whole field, 0.56 among the top fifteen.
%  - Buckley & Voorhees' caveat is about bpref, their own measure, but the assumption is
%    shared with condensed-list measures, which is why the paragraph states it as shared
%    before attributing the caveat.
%  - FAGGIOLI REMOVED 2026-09-24 (Abram), kept in the bib, uncited. ⚠ At 250 citations it
%    is the most-cited paper in this space; the cut is defensible on framing but that is
%    the weight being left out.
%  - SOBOROFF REMOVED 2026-09-24 (Abram). His ceiling objection belongs in §4.4, where the
%    reply is available. Bib entry kept for that use.
%  - ⚠ Abbasiantaeb et al. published a 2025 follow-up partly SOLVING the hole-filling
%    problem with a fine-tuned model, so the 2024 pessimism this section reports is being
%    walked back by its own authors. Not cited; recorded in §16.
%  - ✓ RESOLVED 2026-09-29 (citation audit). froebe2025assessors and
%    szymanski2025judgelimits are fully verified against the published records;
%    upadhyay2024umbrelatool was removed. sources.bib now carries no TODO or UNVERIFIED
%    markers. Neither remaining entry has a PDF in thesis/literature/, but both were
%    checked field by field against Crossref, OpenAlex and the camera-ready.

No collection of any size is judged completely. There are far too many documents, so the
systems taking part submit their rankings and the documents near the top go to the
assessors. Everything else is left unjudged, and unjudged normally counts as not
relevant. That works for the systems that were there at the time. A system built
afterwards retrieves documents nobody ever looked at, and is marked down for finding
things that may have been good answers.

One response is to change the measure rather than the labels. Buckley and Voorhees
proposed bpref, which scores a ranking on its judged documents
alone~\cite{buckley2004incomplete}; Sakai and Kando later deleted judgments from real
collections to see which measures still ordered systems correctly, and found the
condensed family most robust~\cite{sakai2008prime}. Both set aside what nobody looked at,
and both are safe only if the unjudged documents are no better and no worse than the
judged. Buckley and Voorhees decline to conclude that bpref fairly evaluates a system
which did not take part, since a new system can break that
assumption~\cite{buckley2004incomplete}.

The other response is to supply the missing labels. Thomas and colleagues showed that a
carefully written prompt produces labels close to those of trained assessors on web
search, and how delicate that is: across thirty-two variants, agreement ranged from
$\kappa$ 0.20 to 0.62~\cite{thomas2024llm}. An open-source reproduction, UMBRELA, was then run
across a public evaluation of seventy-seven runs from nineteen teams beside the
organisers' own assessors~\cite{upadhyay2025umbrela}. Rankings of
whole systems came out close; agreement topic by topic was weaker, the model was more
lenient, and having humans check its labels did not improve the
match~\cite{upadhyay2025umbrela}.

That result is disputed in three ways. Clarke and Dietz set that overall rank correlation
of 0.89 against the 0.56 they find among the top fifteen systems --- where the
comparison matters~\cite{clarke2024llmassessment}. Fr\"obe and colleagues compared
eight automatic assessors. They agreed with each other more closely than any of them
agreed with a human, and they favoured systems that used a language model to rank,
placing those on average
nine to seventeen positions too high~\cite{froebe2025assessors}. And Abbasiantaeb and
colleagues labelled only the documents nobody had judged, then used those labels beside
the assessors' own; that mixture tracked the assessors worse than labelling every
document automatically did~\cite{abbasiantaeb2024holes}.

All of that evidence comes from web search, general question answering and conversational
search. Where judging needs specialist knowledge the picture is worse: in two
professional fields, subject-matter experts agreed with an automatic judge on roughly two
thirds of cases, 68\% and 64\%~\cite{szymanski2025judgelimits}. Anyone relying on automatic labels
therefore has to show they work on the collection in hand rather than point at results
from another.


\section{Research gap}\label{sec:rw-gap}

% Rewritten 2026-09-24 (Abram), replacing the 2026-09-24 first version. 287 w -> 263 w,
% against a ~280 w budget. Three short paragraphs, one per gap, mapping to RQ1, RQ2 and C3.
% States questions only -- no outcome, no hypothesis ID, no figure of this thesis's own, and
% NO RESEARCH QUESTION NAMED: the reader met those in Chapter 1, and naming them here turns
% the section into a checklist. No new citations either; every source already appears in
% Sections 2.1-2.4.
% The chapter ends here with no forward pointer. Chapter 3 inherits machinery, not an
% argument, so there is no tension for one to resolve.
%
% ⚠ TWO OF THE THREE LEGS WERE FALSE IN THE FIRST VERSION. Both fixed here; do not let
%   either wording return.
%
%  LEG 1 said the formula engines had "never been entered into the answer retrieval task at
%  all, let alone paired with a text retriever". The narrow half survives -- no tree-EMBEDDING
%  engine has been -- but the "let alone" clause is false. Four ARQMath participant systems
%  pair formula structure with text for answer retrieval: MathDowsers (tuples from formula
%  trees indexed beside text terms in one engine, best Task 1 participant run twice, Kane et
%  al. 2022), Approach Zero with Anserini (operator-tree path tokens interpolated with a BM25
%  run, Zhong et al. 2021), Approach Zero with ColBERT (structural matching plus a dense
%  retriever, top across all three ARQMath-3 tasks, Zhong et al. 2022) and Kassaie et al. 2025
%  (RRF over math and keyword components). None is cited in the thesis or in sources.bib.
%  The leg now argues the GROUP claim §2.2 makes: the combining has been done, and in every
%  case the formula side was one index rather than several.
%
%  LEG 2 said "On mathematical questions almost nothing has been measured at all". False:
%  MathNet (2026), Chen et al. 2025 on an undergraduate textbook, Jiang et al. 2026's
%  dual-index framework, Confident RAG (2025) and Arabzadeh et al. 2026 all measure it. See
%  manifest §16's note under §2.4. It also asserted that retrieval quality "turns out to be a
%  poor guide to how good the answer is" -- the flat claim §2.3 retired on 2026-09-24, so the
%  chapter contradicted itself two pages apart. The leg now says the relationship is
%  unsettled, and that what none of that work varies is how formulas are represented.
%
%  LEG 3 SURVIVES AND IS NOW BETTER EVIDENCED. §2.4 carries Szymanski's 68% / 64% expert
%  agreement, so this leg leans on a measured finding instead of asserting. Of the three, it
%  is the only one the 2026-09-24 searches strengthened.
%
% Also fixed: the first version re-narrated Henkel, a study the reader met in §2.3, and
% repeated its overstated null. A gap section names what is missing; the evidence lives in
% §2.1-§2.4. And all three paragraphs closed on the same frame -- "is not known", "has not
% been asked", "has not been established". The closers now differ.
%
% The shape each leg follows, and the reason it is repeated: state what HAS been done, then
% the one variation this thesis makes that has not. An absence claim invites a counterexample;
% two of them got one. Naming the adjacent work first is what makes the remaining gap hard to
% attack.

The first gap is on the retrieval side. Formula representations are built and tuned for
formula search, and that is where they are compared against one another. Systems built to
retrieve answers have combined formulas with text before. In each of them the formulas
were handled by a single index, not several. The only system that keeps several formula
representations apart does so in formula search, where there is no text to combine them
with~\cite{mansouri2022dprl}. So the grouping of those representations, and the way that
group meets text retrieval, is untried where answers are what is retrieved.

The second gap is at the far end of the pipeline. Whether better retrieval yields better
answers is unsettled in general, and on mathematics it has been measured only over curated
corpora of prepared material, textbook sections and problem sets retrieved by semantic
similarity~\cite{henkel2024mathrag,alshammari2026mathnet}. None of that work varies how
formulas are represented in the retriever. So whether a difference that is visible at the
ranking survives into the answer, on mathematics, has not been asked.

The third gap is in whether any of it can be measured. A system built after a collection
was judged is scored against judgments made for other systems, and the two remedies,
changing the measure or supplying the missing labels, both rest on assumptions the
literature has not settled. Automatic labels have been tested on web search, on general
question answering and on conversational search. Where judging needs specialist knowledge,
they agree with human experts less often than those experts agree with one
another~\cite{szymanski2025judgelimits}. No one has
tested them on mathematical answers, where relevance turns on whether an explanation
actually helps.
